\subsection{THETA3-C25: a concrete exceptional Weil source with endomorphism ring $\Z[i]$}\label{C25}

The preceding non-CM sources can be sharpened to a geometrically
simple Weil-type fourfold with no extra endomorphisms. Its Weil
space is entirely outside the divisor algebra. Thus the
projector criterion C13 now has a specified exceptional target
in the same arithmetic ball covering tower.

\paragraph{The algebraic torsion prescription.}
Keep
\[
 E:y^2=x^3-30x-25,\qquad T=(-5,0).
\]
Fix an embedding $\overline\Q\hookrightarrow\overline\Q_{109}$.
Write $s^2=2$ in $\mathbf F_{109^2}$; two is a nonsquare
modulo 109. Let $U,V,W\in E[8](\overline\Q)$ and
$R\in E[13](\overline\Q)$ be the unique prime-to-109
torsion lifts of the following points:
\[
\begin{array}{c|cc}
 &x&y\\\hline
 U&34+53s&11+96s\\
 V&2+18s&8+82s\\
 W&2+91s&8+27s\\
 R&47&20s .
\end{array}
 \tag{C25a}
\]
Here uniqueness is in the indicated finite \'etale torsion
group over the unramified quadratic extension of $\Q_{109}$.
These are algebraic points, so this prescription defines
objects over a number field, not merely over a $p$-adic field.
For $\mathcal P=\{U,U+R,V,W\}$ put
\[
 H:\quad y_H^2=F_H(x),\qquad
 F_H(x)=\bigl((x+5)^2-45\bigr)
       \prod_{P\in\mathcal P}(x-x(P))(x-x(P+T)),
 \tag{C25b}
\]
and take its smooth projective model. Then $H$ has genus four,
and its Jacobian $A$ satisfies
\[
 \operatorname{End}_{\overline\Q}(A)=\Z[i],\qquad
 \operatorname{NS}(A)_\Q=\Q\Theta.
 \tag{C25c}
\]
It is geometrically simple and is of $\Q(i)$-Weil type with
signature $(2,2)$. For a connected finite \'etale cover
$X\to Y$ of Schoen's ball surface one has
\[
 \deg(X/Y)\leq43264,\qquad
 \operatorname{Alb}(X)\twoheadrightarrow A.
 \tag{C25d}
\]
Its two-dimensional Weil space $W_{\Q(i)}(A)$ meets the
subspace generated by products of rational divisor classes
only in zero.

\paragraph{The torsion calculations and the good reduction.}
All the following calculations use the displayed addition law
of $E$ and arithmetic in $\mathbf F_{109}[s]/(s^2-2)$.
They are supplied with independent exact certificates.
The four points in (C25a) lie on $E$, and
\[
 2U=(34,0),\quad4U=0,\qquad
 4V=4W=T,\quad8V=8W=0,\qquad13R=0,
\]
with $R\ne0$. If $\varphi$ denotes 109-power Frobenius,
\[
 \varphi(U)=-U+T,\quad
 \varphi(V)=W,\quad\varphi(W)=V,\quad
 \varphi(R)=-R.
 \tag{C25e}
\]
The three four-element orbits of $U,V,W$ under negation
and translation by $T$ are distinct and avoid both
$E[2]$ and the solutions of $2P=T$. Injectivity of
prime-to-109 reduction proves the same assertions for
their lifts. The order-thirteen point cannot identify
the orbit of $U+R$ with any of these order-eight orbits.
Thus (C25b) has ten distinct branch points, all in the
image of $E[104]$ on its Kummer line. The even-level
transfer of C24 proves (C25d), since $4\cdot104^2=43264$.
The coordinate change (C24c) converts the branch divisor
to $0,\infty$ and four opposite pairs. Hence the balanced
$\Q(i)$-action and its Rosati property follow from the
explicit calculation in C24.

The Frobenius relations (C25e) hold also for the unique
torsion lifts over the unramified quadratic extension.
They preserve the entire branch divisor and its monic
equation. Thus the coefficients of (C25b) lie in
$\Q_{109}$ under the chosen embedding, and the number
field generated by those coefficients has a place with
residue field $\mathbf F_{109}$. This is stronger than
merely observing that a polynomial over a larger residue
field happens to have coefficients in the prime field.
At this place its reduction is
\[
\begin{split}
 \overline F_H(x)={}&x^{10}+62x^9+27x^8+30x^7+45x^6
       +3x^5+40x^4\\
       &+35x^3+65x^2+x+63 .
\end{split}
 \tag{C25f}
\]
Euclid's algorithm gives
$\gcd(\overline F_H,\overline F_H')=1$.
Its leading coefficient is one, so this is a smooth
genus-four reduction with two points at infinity.

\paragraph{A finite absolute-simplicity certificate.}
The exact counts on this one curve are
\[
\begin{array}{c|rrrr}
 r&1&2&3&4\\\hline
 \#H(\mathbf F_{109^r})&134&11794&1294958&141166026 .
\end{array}
 \tag{C25g}
\]
The accompanying computations enumerate all elements of
the finite fields and use quadratic characters to count
the two-sheeted cover, including its two points at infinity.
One implementation uses
$\mathbf F_{109}[u]/(u^2-2)$ and
$\mathbf F_{109}[w]/(w^4-2)$; the independent implementation
uses $u^2-6$ and $w^4-6$ instead.
For the quartic fields irreducibility follows by taking
the quadratic tower: a nonsquare $d$ has square root
$u$ whose norm is $-d$, a nonsquare because
$109\equiv1\pmod4$; thus $u$ is a nonsquare in
$\mathbf F_{109^2}$.
For each element $f$ the identity
\[
 \chi_{\mathbf F_{109^r}}(f)
  =\chi_{\mathbf F_{109}}
       \bigl(N_{\mathbf F_{109^r}/\mathbf F_{109}}f\bigr)
\]
is used, with both sides zero at zero. It follows from
the geometric-series formula for the norm exponent.
All enumerations and norm operations are exact integer
arithmetic modulo 109.

Put $s_r=109^r+1-\#H(\mathbf F_{109^r})$.
Newton's recursion
$k a_k=-\sum_{j=1}^k a_{k-j}s_j$, $a_0=1$,
followed by the reciprocal functional equation, gives
the Frobenius polynomial
\[
\begin{split}
 f(Z)={}&Z^8+24Z^7+244Z^6+1224Z^5+3510Z^4\\
       &+133416Z^3+2898964Z^2+31080696Z+141158161.
\end{split}
 \tag{C25h}
\]
The middle coefficient is a unit modulo 109, so the
reduction is ordinary. Here is an exact irreducibility
and geometric-simplicity test, rather than a numerical
root comparison.

Over $\Q(i)$ one has $f=g\overline g$, where
\[
 g(Z)=Z^4+(12+10i)Z^3+192iZ^2+(-1308+1090i)Z-11881.
 \tag{C25i}
\]
Modulo the Gaussian prime defined by $61=0$, $i=11$,
this is $Z^4-23Z^2+7Z+14$.
In its polynomial quotient over $\mathbf F_{61}$,
\[
 Z^{61^4}=Z,\qquad
 \gcd(Z^4-23Z^2+7Z+14,Z^{61^2}-Z)=1.
 \tag{C25j}
\]
These modular-power and Euclidean checks prove its
irreducibility. Hence $g$ is irreducible over $\Q(i)$.
If $\pi$ is one of its roots, its minimal polynomial over
$\Q$ is divisible by $g$ after extension to $\Q(i)$;
its rational coefficients make it divisible by
$\overline g$ as well. These two irreducible quartics are
distinct, so their product $f$ is irreducible over $\Q$.
The reduction of $A$ is consequently simple.

Write $K_\pi=\Q[Z]/(f)$ and again write $\pi$ for the
class of $Z$. We verify that $\Q(\pi^d)=K_\pi$ for
every $d>0$ by Howe--Zhu, \emph{On the existence of
absolutely simple abelian varieties of a given dimension
over an arbitrary field}, Proposition~3 and Lemma~5,
preprint pp.~4--5. Their criterion says that a first
failure would either make $f$ a polynomial in $Z^d$,
or would put a primitive $d$th root of unity in
$K_\pi$ with $\Q(\pi^d)$ a proper subfield.
The nonzero $Z^7$ coefficient rules out the first case.
In the second case $\phi(d)$ divides $[K_\pi:\Q]=8$,
so the only $d>1$ to test are
\[
 2,3,4,5,6,8,10,12,15,16,20,24,30.
 \tag{C25k}
\]
This list follows directly from
$\phi(\prod\ell^e)=\prod\ell^{e-1}(\ell-1)$:
the only possible prime factors are $2,3,5$, and their
allowed powers give precisely the list displayed.

For each $d$ in this order, express
$1,\pi^d,\ldots,\pi^{7d}$ in the basis
$1,\pi,\ldots,\pi^7$. The determinants of their
integer coefficient matrices, reduced modulo 37, are
\[
 5,11,22,30,9,18,28,14,20,33,7,17,22.
 \tag{C25l}
\]
All are nonzero. Hence each is a rational basis of
$K_\pi$, ruling out all the possible failures in
(C25k). Howe--Zhu's proposition now proves that the
reduction is absolutely simple. The independent
certificate reproduces (C25i)--(C25l), including the
power-basis determinants, by modular polynomial
arithmetic and Gaussian elimination.

The generic $A$ is geometrically simple as well.
Otherwise complete reducibility supplies a nontrivial
idempotent in its geometric rational endomorphism
algebra. After a finite extension it specializes
injectively at the good place to the rational
endomorphism algebra of the reduction; see Conrad,
\emph{Gross--Zagier revisited}, Theorem~2.1(2),
printed p.~70. This contradicts the absolute simplicity
of that reduction, whose rational endomorphism algebra
is a division algebra. Thus no such idempotent exists.

\paragraph{A rank-two torus forces the exact endomorphism ring.}
We now verify that the very same curve has the two-node
reduction used in C24 at a place above thirteen.
All nonzero points of $E[13]$ are conjugate over $\Q$.
Indeed $E(\mathbf F_7)$ has six points, so Frobenius at
seven has characteristic polynomial
$Z^2-2Z+7$ on $E[13]$.
It is irreducible over $\mathbf F_{13}$.
If $\beta$ is its root in $\mathbf F_{169}$, direct
modular powering gives
\[
 \beta^{168}=1,\quad
 \beta^{84}=-1,\quad
 \beta^{56}=9,\quad
 \beta^{24}=2+2\beta.
 \tag{C25m}
\]
The last three values are not one; since $2,3,7$ are
the prime divisors of 168, $\beta$ has order 168.
Multiplication by it acts transitively on the 168
nonzero elements of the two-dimensional $\mathbf F_{13}$
space. The global Galois action contains this Frobenius
action, proving the assertion.

The displayed model of $E$ has good reduction at thirteen.
As in C24, reduction of its 169 geometric thirteen-torsion
points has a nonzero kernel, since its special elliptic
curve has at most thirteen geometric thirteen-torsion
points. Transitivity therefore gives a place above
thirteen at which the particular $R$ in (C25a) reduces
to zero. Reduction remains injective on $E[8]$ at this
place. Apply exactly the two-node argument of C24,
with the good prime thirteen in place of seven:
the branch pairs from $U$ and $U+R$ coalesce, the
other branch values stay distinct, and the geometric
stable fibre is irreducible with two nodes. Hence
$A$ has semistable toric rank exactly two, persistent
under every further finite extension.

Let $D=\operatorname{End}_{\overline\Q}^0(A)$.
Geometric simplicity makes $D$ a division algebra.
After defining all endomorphisms over one finite
extension, the N\'eron property lets them act on the
two-dimensional geometric torus in this semistable
fibre. Its rational character space has dimension two,
so one obtains a unital algebra homomorphism
\[
                    D\longrightarrow M_2(\Q).
\]
It is injective: the kernel is a two-sided ideal of a
division algebra and the identity acts nontrivially.
Equivalently, a nonzero endomorphism of the simple
$A$ is an isogeny, and an inverse up to multiplication
by an integer induces an inverse on the torus up to
the same integer. The two-dimensional rational
character space is a nonzero vector space over the
division algebra $D$ (using the opposite algebra if
pullback conventions require it). Therefore
$\dim_\Q D$ divides two.
The existing $\Q(i)$-action forces $D=\Q(i)$.
Every integral endomorphism is an algebraic integer,
and the curve automorphism already supplies $\Z[i]$.
As this is the maximal order in $\Q(i)$, the integral
endomorphism ring is exactly $\Z[i]$.

The polarization identifies rational divisor classes
with Rosati-symmetric rational endomorphisms: Milne,
\emph{Abelian Varieties}, version~2.00 (2008),
Proposition~14.2, printed pp.~61--62. The
Rosati involution on $\Q(i)$ is conjugation, whose
fixed field is $\Q$, proving the N\'eron--Severi
assertion in (C25c). Alternatively this identification
is the standard map $L\mapsto\lambda_\Theta^{-1}\phi_L$;
its kernel is $\operatorname{Pic}^0$ and its image is
the symmetric subspace.
Every codimension-two divisor polynomial is thus a
multiple of $\Theta^2$. Under the two $K$-eigenspaces
of $H^1$, the polarization has bidegree $(1,1)$ and
$\Theta^2$ has bidegree $(2,2)$. The two Weil lines
have bidegrees $(4,0)$ and $(0,4)$. These are different
summands of $\bigwedge^4H^1(A,\C)$, so their intersection
with the divisor algebra is zero. \hfill$\square$

\paragraph{The remaining cycle test is now on a named exceptional target.}
Apply C13 to the torsion-prescribed fourfold (C25b).
One nonzero projected surface class spans all of its
Weil space under $\Q(i)$. C14 proves that surfaces
whose maps to $A$ factor through products of two
curves have zero such projection. Thus the actual
non-CM source, its simplicity, its endomorphism
algebra, and the exceptional nature of the target
space have all been decided here. The new task is
the specific projected correspondence, with these
obstructions already accounted for. The supplied
Markman theorem already proves algebraicity of
Weil classes on fourfolds; this finding is not a
new proof of that theorem. It gives an explicit
exceptional test object accessible to the ball
construction, and does not assert the Hodge
conjecture for higher-dimensional non-CM targets.

\framework{The source correspondence and Weil projector are in
the homological-motive framework of cote~44, scans~7--8,
and cote~16, scans~43--46. The persistent torus and
its endomorphism action are the N\'eron--Picard and
dual-graph mechanisms of SGA~7~I, IX~12.1(d),
printed pp.~467--468, and IX~12.3.6--7,
printed pp.~471--472. The finite-field calculations
certify a specified source and its endomorphisms;
they do not lift a vertical arithmetic intersection
to an unconstructed characteristic-zero cycle.}
